\documentclass[11pt]{article}
\usepackage[margin=1in]{geometry}
\usepackage{amsmath,amssymb}
\usepackage{enumitem}
\usepackage{microtype}
\usepackage[hidelinks]{hyperref}

\setlength{\parindent}{0pt}
\setlength{\parskip}{4pt}

\newlist{qna}{enumerate}{1}
\setlist[qna]{leftmargin=1.6em,itemsep=3pt,topsep=4pt,label=\textbf{\arabic*.}}
\newcommand{\paper}[1]{\subsection*{#1}}

\newcommand{\q}[2]{\item \textbf{#1} #2}
\newcommand{\tiein}[1]{\par\smallskip\noindent\emph{Connection to the offline RL problem space and BCA:} #1\par\medskip}

\title{Literature Review: Offline Reinforcement Learning\\Parts A and B}
\author{David Bayha\\ITSC 8110}
\date{October 2026}

\begin{document}
\maketitle



\section*{Part A: Research Papers}



% ============================ PAPER 1 ============================
\paper{Paper 1: Conservative Q-Learning for Offline Reinforcement Learning (CQL)}
CQL trains the critic to be deliberately pessimistic about any action the dataset does not contain, so the policy stops chasing accidentally inflated scores.
\begin{qna}

\q{What is the citation?}
Kumar, A., Zhou, A., Tucker, G., \& Levine, S. (2020). Conservative Q-learning for offline reinforcement learning. \emph{Advances in Neural Information Processing Systems, 33}, 1179--1191. Version reviewed: arXiv:2006.04779v3. \url{https://doi.org/10.48550/arXiv.2006.04779}

\q{What is the research area?}
Machine learning, specifically offline reinforcement learning (RL): learning a policy from a fixed log of past decisions with no further trial and error. Within offline RL, CQL belongs to the ``pessimistic critic'' family: instead of restricting the policy, it lowers the scores assigned to actions outside the data. This distinguishes it from policy-constraint methods (BEAR, BRAC) and ensemble/uncertainty methods.

\q{What is the research question?}
Can a critic trained to be systematically pessimistic about unseen actions produce a policy whose predicted value is a guaranteed floor on its true performance, and does that beat methods that instead constrain the policy to stay near the recorded behavior?

\q{What is the significance of the question?}
Offline RL is most valuable exactly where live experimentation is expensive or dangerous (robotics, medicine, recommendation) yet where large logs already exist. Standard algorithms fail in this setting because the critic's guesses about unseen actions are never corrected and the policy is drawn to precisely those inflated scores; earlier fixes (modeling the behavior policy, uncertainty penalties) proved fragile. A cheap, principled fix would let RL reuse logged data the way supervised learning reuses datasets.

\q{What are the contributions reported in this paper?}
The novel idea is a simple addition to the critic's training objective: push down the scores of whatever actions the critic currently likes too much, while propping up the scores of actions that actually appear in the data, so the net penalty lands mainly on unseen actions. Around this the paper provides two ready-to-use variants (CQL($\mathcal{H}$) and CQL($\rho$)), mathematical guarantees that the learned value undercuts the true value and that the policy safely improves on the one that collected the data, and a roughly twenty-line modification of two standard algorithms, with an automatic scheme for setting the caution strength $\alpha$.

\q{What is the research method?}
A mix of algorithm design, mathematical analysis (mostly for simplified tabular cases, with extensions), and large-scale benchmarking on the D4RL suite (simulated locomotion, mazes, dexterous-hand and kitchen tasks) and offline Atari against established prior methods. Diagnostics check directly whether predicted values stay below true returns, and ablations isolate key design choices.

\q{What are the results?}
On datasets recorded by a single behavior, CQL is close to the best prior method; on the more realistic mixed-source datasets it usually wins by large margins (e.g., 91.1--110.5 on the three random-expert locomotion tasks where the best competitors score 53.0, 2.7, and 11.1), and only the CQL variants score above zero on the hardest mazes. As designed, its predicted values are confirmed to sit below true returns, while competing methods overestimate.

\q{What are the impacts and limitations of the contributions?}
CQL became the field's standard pessimistic baseline and showed pessimism can be added cheaply, without modeling the behavior that produced the data. Its limits: the guarantees cover simplified settings rather than the deep networks actually used; the caution knob $\alpha$ (or its auto-tuning threshold) must be retuned per domain; there is no way to decide when to stop training without trying the policy live; and one global $\alpha$ treats well-covered and poorly covered situations identically.

\q{How were the results evaluated or validated?}
By mathematical proofs in appendices, D4RL scores averaged over four runs, direct measurements comparing predicted values against actual returns, and ablations. The evidence is broad but statistically thin: means without spreads, some baselines copied rather than rerun, and a few inconsistent tables.

\q{What are the ethical implications? How are they addressed?}
The reviewed version contains no broader-impact discussion. The ``safe improvement'' guarantee is only relative to the logging policy, probabilistic, and assumption-bound, so the word \emph{safe} invites misplaced trust; a policy kept near the data also keeps the data's biases; all validation is simulated; compute cost is unreported. Conservatism is an indirect step toward safer learning, but bias, deployment safeguards, and societal consequences go undiscussed.

\q{What themes are suggested by this research? (minimum 3)}
\emph{Distribution shift}: scores for unseen actions cannot be corrected offline. \emph{The conservatism--improvement trade-off}: pessimism must be dosed, here via $\alpha$. \emph{Critic-side vs.\ policy-side restraint}: shown to be two faces of one idea. \emph{Theory vs.\ practice}: the guarantees assume what deep networks violate. \emph{Dataset quality}: mixed data is where caution pays.

\end{qna}
\tiein{CQL defines the ``one global pessimism knob'' paradigm of the problem space; applied to CQL, BCA (Bayesian Conformal Aggregation) keeps the critic-side idea but adds calibrated per-sample caution on top of the single $\alpha$, a move CQL's own appendix hints at with a count-dependent penalty.}


% ============================ PAPER 2 ============================
\paper{Paper 2: A Minimalist Approach to Offline Reinforcement Learning (TD3+BC)}
TD3+BC shows that adding one ``stay near the recorded actions'' penalty to a standard online algorithm matches far more complicated offline methods at a fraction of the compute.
\begin{qna}

\q{What is the citation?}
Fujimoto, S., \& Gu, S. S. (2021). A minimalist approach to offline reinforcement learning. \emph{Advances in Neural Information Processing Systems, 34}, 20132--20145. Version reviewed: arXiv:2106.06860v2. \url{https://doi.org/10.48550/arXiv.2106.06860}

\q{What is the research area?}
Offline RL for continuous control (robot-like tasks with smooth action spaces). TD3+BC belongs to the behavior-regularization family: methods that fight inflated scores for unseen actions by keeping the policy close to the recorded behavior.

\q{What is the research question?}
Can a standard online algorithm, changed as little as possible, reach state-of-the-art offline performance, and how much of the extra machinery in recent offline methods is actually necessary?

\q{What is the significance of the question?}
Every extra component (a learned model of the recorded behavior, pre-training phases, architecture changes) is a choice that normally must be validated by trying policies, which requires the very interaction offline RL is supposed to avoid. A minimal method isolates what actually works, raises the bar that elaborate methods must clear, and stays cheap as datasets grow.

\q{What are the contributions reported in this paper?}
Three. First, a detective-style audit showing that two leading methods (CQL and Fisher-BRC) owed much of their scores to poorly justified engineering tweaks, not just their stated ideas. Second, the TD3+BC algorithm itself: one penalty term in the actor's objective (``score high, but stay close to the recorded action'') plus standardized inputs, with a single balance setting reused across all tasks. Third, a model of careful practice: unified reruns of competitors, wall-clock timings, a demonstration that offline policies are unstable across episodes, and public code.

\q{What is the research method?}
The study combines critical analysis with careful experimentation rather than new mathematics. The actor maximizes $\lambda Q(s,\pi(s)) - (\pi(s)-a)^2$ averaged over recorded state-action pairs, with the balance weight $\lambda$ automatically rescaled by the typical size of the scores so one value ($\alpha = 2.5$) works everywhere. Experiments cover fifteen standard locomotion datasets with five random seeds, plus ablations removing each ingredient.

\q{What are the results?}
Summed over the fifteen tasks, TD3+BC essentially ties the best competitor (979.3 vs.\ 974.6) and beats CQL by a wide margin (764.3), while training in 39 minutes versus over four hours for CQL. Ablations confirm both the RL term and the penalty are needed; normalization gives a smaller boost. It fails where imitation fails: datasets heavy in poor behavior (28.8 vs.\ CQL's 100.3 on a half-random, half-expert Walker2d mix) and the hardest mazes.

\q{What are the impacts and limitations of the contributions?}
TD3+BC became the field's default baseline and the yardstick new methods must beat. It also urged the community to report stability, not just averages. Limitations: no theory; a uniform pull toward recorded actions cannot tell good recordings from bad ones and can average over genuinely different behaviors; the diagnosed instability is left unsolved; and every sample receives the same constraint strength regardless of how well the data cover it.

\q{How were the results evaluated or validated?}
Five seeds, ten test episodes per checkpoint, scores averaged over the final checkpoints rather than the luckiest one, and the two strongest competitors (CQL and Fisher-BRC) rerun from their own authors' code under one protocol. Weaknesses: no statistical tests or confidence intervals, only five seeds despite high documented variance, and a tie with Fisher-BRC that supports parity rather than superiority.

\q{What are the ethical implications? How are they addressed?}
A short statement notes that learning from logged data inherits its biases, warns of misuse (e.g., more habit-forming recommendation), and argues the main effect is wider access. The paper's most useful safety contribution is candid documentation that offline policies can fail badly on individual episodes. No mitigations are proposed; all evidence is simulated.

\q{What themes are suggested by this research? (minimum 3)}
\emph{Minimalism}: a one-line change suffices. \emph{Behavior regularization}: the uniform pull toward the data. \emph{Hidden implementation details}: engineering, not just ideas, drove prior scores. \emph{The difficulty of truly offline tuning}: validation needs the interaction offline RL avoids. \emph{Evaluation rigor and reproducibility}; \emph{computational efficiency}.

\end{qna}
\tiein{TD3+BC is the minimalist pole of the problem space, with one uniform constraint for every sample. Its successor (ReBRAC, Paper 4) showed this constraint is what carries the method; BCA strengthens exactly this term per sample by a calibrated factor, keeping the global weight.}

% ============================ PAPER 3 ============================
\paper{Paper 3: Offline Reinforcement Learning with Implicit Q-Learning (IQL)}
IQL learns values using only the actions already in the dataset, never asking ``what if'' about unseen actions, yet can still stitch together good behavior from suboptimal pieces.
\begin{qna}

\q{What is the citation?}
Kostrikov, I., Nair, A., \& Levine, S. (2022). Offline reinforcement learning with implicit Q-learning. In \emph{International Conference on Learning Representations (ICLR 2022)}. Version reviewed: arXiv:2110.06169v1. \url{https://doi.org/10.48550/arXiv.2110.06169}

\q{What is the research area?}
Offline RL. IQL takes a third path beside constraining the policy (TD3+BC) and penalizing the critic (CQL): \emph{in-sample} learning, in which value networks are only ever evaluated on recorded actions, and the policy is extracted afterward by imitation weighted toward better-than-expected actions (AWR).

\q{What is the research question?}
Can an offline algorithm still plan ahead by chaining many steps of value backups toward optimal behavior while its value functions are trained only on logged actions, so that unseen actions are never scored at all?

\q{What is the significance of the question?}
Asking the critic about unseen actions is the root cause of offline failure, and every existing fix compromises between safety and improvement. One-step in-sample alternatives avoid the problem but cannot stitch fragments of good behavior together; for example, they cannot assemble a full maze solution from partial trajectories. Removing the problem while keeping multi-step planning would resolve the field's central tension, cheaply.

\q{What are the contributions reported in this paper?}
The novel core is ``implicit'' value learning: a state-value network trained with an asymmetric loss (an expectile loss, whose level acts as an optimism dial) estimates the best value achievable within the data, and a second network then learns action values one step further, so the pair plans ahead without ever scoring an unseen action. Supporting contributions: a proof that turning the dial fully optimistic recovers the best in-data value; policy extraction by weighted imitation; a correction of earlier unfair benchmark comparisons; an offline-to-online fine-tuning study; and fast public code.

\q{What is the research method?}
Algorithm design plus a short limiting argument, then benchmarking: a toy maze built to require stitching; standard locomotion, maze, hand-manipulation, and kitchen suites against up to seven baselines; an ablation of the optimism dial; and a fine-tuning study in which offline training is followed by live interaction.

\q{What are the results?}
On locomotion, IQL roughly ties the best prior method (692.4 vs.\ 698.5); on the maze tasks that demand stitching it dominates (378.0 vs.\ 303.6 total; 39.6 and 47.5 on the largest mazes, where the strongest competitor manages about 15). Training takes about 20 minutes. After live fine-tuning it improves further (370 to 474 on mazes) where CQL collapses on some tasks. It still loses individual tasks, notably to filtered imitation on one locomotion dataset.

\q{What are the impacts and limitations of the contributions?}
IQL showed in-sample learning need not give up long-horizon planning, and because its value learning never depends on the policy, the extraction step can be swapped or continued online. Limitations: its two knobs are set per domain; the tables report means without spreads; continuous states usually carry only one recorded action, so the ``best in-data value'' relies on the network generalizing across similar states; and its fixed weighting amplifies trustworthy and noisy advantages identically.

\q{How were the results evaluated or validated?}
Standard benchmark averages (ten seeds on locomotion, 100 episodes on mazes), a corrected all-one-version comparison, and one ablation. Weaknesses: baseline numbers from mixed sources (one competitor's maze scores could not be reproduced and were copied), no reported spreads, no stated procedure for choosing per-domain settings, and only the optimism dial ablated.

\q{What are the ethical implications? How are they addressed?}
No ethics discussion. Because the policy is built by reweighting recorded behavior, biases in the log are inherited and can be amplified whenever they scored well; the deliberate optimism of the upper expectile is risky in safety-critical use; and the fine-tuning recipe involves live exploration that would carry physical risk on hardware. Low compute cost is a modest positive.

\q{What themes are suggested by this research? (minimum 3)}
\emph{Avoiding distribution shift by design}: never query unseen actions. \emph{The conservatism--improvement trade-off as a single dial}: the expectile level. \emph{Constraint relocated}: the imitation weights, not the critic, hold the caution. \emph{Dataset composition decides which family wins}: locomotion lets one-step methods compete; mazes demand stitching. \emph{Minimalism}; \emph{offline-to-online fine-tuning}.

\end{qna}
\tiein{IQL is one of BCA's four host algorithms: its imitation weight is exactly where calibrated per-sample widths plug in, tempering inflated advantages while keeping IQL's one global strength.}

% ============================ PAPER 4 ============================
\paper{Paper 4: Revisiting the Minimalist Approach to Offline Reinforcement Learning (ReBRAC)}
ReBRAC shows that giving the simple TD3+BC baseline the same quiet engineering upgrades newer methods enjoyed, including deeper networks, normalization, and separate penalties, is enough to match or beat far more complex algorithms.
\begin{qna}

\q{What is the citation?}
Tarasov, D., Kurenkov, V., Nikulin, A., \& Kolesnikov, S. (2023). Revisiting the minimalist approach to offline reinforcement learning. \emph{Advances in Neural Information Processing Systems, 36}, 11592--11620. Version reviewed: arXiv:2305.09836v2. \url{https://doi.org/10.48550/arXiv.2305.09836}

\q{What is the research area?}
Offline RL, in the behavior-regularization lineage (BRAC, TD3+BC). Deliberately not a new principle: it is a systematic study of which ``incidental'' design choices bundled with new algorithms actually improve an established baseline.

\q{What is the research question?}
How far can a minimalist baseline be improved purely by adopting the small design choices that recent methods shipped quietly alongside their headline ideas (deeper networks, layer normalization, bigger batches, separate penalty weights for actor and critic, and a longer planning horizon), and which of those choices account for the gain?

\q{What is the significance of the question?}
If gains credited to new algorithms actually come from quiet engineering, the field mis-attributes progress and practitioners cannot tell what is worth adopting. Establishing how strong a cheap baseline can become resets the bar for everyone and gives practitioners a fast default. The elaborate alternatives were expensive (one needed as many as 500 critic networks).

\q{What are the contributions reported in this paper?}
An audit cataloging the hidden engineering inside seven recent methods; ReBRAC itself (TD3+BC plus five upgrades, including separate ``stay near the data'' penalties on the actor and on the critic); a broad 51-dataset evaluation spanning state-based and image-based benchmarks plus an offline-to-online study; careful ablations identifying what carries the gains; and a full reproducibility package (code, per-dataset settings, compute tables).

\q{What is the research method?}
Design by composition, then large-scale benchmarking and ablation; no new theory. Settings are tuned on four random seeds, and the chosen configuration is then re-evaluated on ten fresh seeds using final-checkpoint reporting, a fairer protocol than reporting lucky peaks. Diagnostic runs remove or add one upgrade at a time.

\q{What are the results?}
ReBRAC is the best ensemble-free method on average on mazes (76.8) and hand-manipulation tasks (58.6), a close second on locomotion (81.2 vs.\ SAC-RND's 82.6), and fine-tunes online from 55.0 to 91.1 with the lowest regret of the methods tested. The ablations are the real finding: remove the actor's stay-near-the-data penalty and the method collapses (72.2 to 11.4); remove critic layer normalization and the score halves; the critic-side penalty adds almost nothing. No single upgrade alone rescues plain TD3+BC.

\q{What are the impacts and limitations of the contributions?}
ReBRAC reset the field's baseline and gave practical guidance: the actor penalty and critic normalization are essential, and the planning horizon matters when rewards are sparse. Limitations: its two penalty weights are still re-tuned per dataset by trying policies live (best values span two orders of magnitude); most tasks recorded by humans remain unsolved; there is no theory; everything is simulated; and constraint strength never varies \emph{within} a dataset.

\q{How were the results evaluated or validated?}
More rigorously than usual: tuning seeds separated from reporting seeds, final-checkpoint reporting, modern aggregate statistics, and a control showing the gains come from the design choices rather than implementation differences. Still, model selection uses real environment returns, which a true offline user lacks, and copied baselines come from weaker protocols.

\q{What are the ethical implications? How are they addressed?}
Not discussed beyond scientific scope. Implicitly: staying near logged behavior inherits its biases and unsafe habits; tuning by live trial runs contradicts the offline premise in risky domains; and performance dipped on one task during live fine-tuning, which outside simulation would mean worse behavior in operation. Compute (roughly 18{,}500 GPU-hours) is disclosed but not discussed.

\q{What themes are suggested by this research? (minimum 3)}
\emph{Implementation details drive progress}: incidental choices explain much ``algorithmic'' advance. \emph{Minimalism strengthened before complexity is added}. \emph{Behavior regularization}: the actor penalty is the indispensable safeguard. \emph{Evaluation rigor}: re-tuning baselines under the same protocol. \emph{Online tuning vs.\ offline model selection}; \emph{offline-to-online fine-tuning}.

\end{qna}
\tiein{ReBRAC's ablations identify the actor's stay-near-the-data term as the lever that matters, but its strength remains a single scalar re-tuned per dataset by live evaluation. BCA's calibrated per-sample weighting is the natural next step: same lever, dosed per sample on top of the host's tuned scalar.}


% ============================ PAPER 5 ============================
\paper{Paper 5: Importance Weighted Actor-Critic for Optimal Conservative Offline Reinforcement Learning (A-Crab)}
A-Crab proves that a small change in how the critic's errors are penalized, by weighting each sample's error instead of squaring it, turns a slow-but-practical method into one with the best possible statistical guarantee.
\begin{qna}

\q{What is the citation?}
Zhu, H., Rashidinejad, P., \& Jiao, J. (2023). Importance weighted actor-critic for optimal conservative offline reinforcement learning. \emph{Advances in Neural Information Processing Systems, 36}, 49579--49602. Version reviewed: arXiv:2301.12714v2. \url{https://doi.org/10.48550/arXiv.2301.12714}

\q{What is the research area?}
Theoretical offline RL with general function approximation, meaning its guarantees are not tied to tiny toy problems. The work connects two lines: adversarial actor-critic methods (ATAC), in which the critic plays devil's advocate against the policy, and importance weighting, which corrects for how rarely parts of behavior appear in the data.

\q{What is the research question?}
Can one practical algorithm achieve the statistically optimal learning speed (error shrinking like $1/\sqrt{N}$ as the data grow) under a weak, averaged assumption about how well the data cover alternative behaviors, where prior practical methods were provably slower or needed unrealistically strong coverage?

\q{What is the significance of the question?}
A guarantee is only as useful as its assumptions. Real logs (medical records, robot traces) never cover every possible policy, and worst-case coverage assumptions collapse when a single rarely logged action makes a likelihood ratio enormous. A practical method with the optimal rate under weak, averaged coverage would narrow the gap between offline RL theory and practice.

\q{What are the contributions reported in this paper?}
The novel step is replacing the squared Bellman-error penalty (a Bellman error is how much the critic's predictions disagree with their own one-step lookahead) with the largest \emph{importance-weighted average} error, where the weights emphasize the samples that matter for the policies being compared. Around this: an optimal-rate proof that drops the strong ``completeness'' closure assumption predecessors needed; a weaker coverage measure; a guarantee of improvement over the logging policy across a wide range of penalty weights; a proof that the old squared penalty fails to reach the optimal rate on some problems; and a practical variant with public code.

\q{What is the research method?}
Mathematics first: a game in which the actor tries to beat the logged behavior while the critic minimizes that improvement plus the weighted error penalty, analyzed with finite-sample concentration tools. A tractable approximation then restricts the weights to a boxed range, yielding a closed-form penalty, and is benchmarked against ATAC on twelve simulated locomotion datasets.

\q{What are the results?}
In theory, A-Crab attains the optimal $1/\sqrt{N}$ rate against any comparator policy the data cover, improving on its predecessors' slower rates ($1/N^{1/3}$ and $1/N^{1/6}$), and improvement over the logging policy is guaranteed across a wide range of penalty settings. In practice it usually trains more stably than ATAC and scores better on several datasets (e.g., about 110 vs.\ a noisy 90--100 on one locomotion task), though results appear only as training curves.

\q{What are the impacts and limitations of the contributions?}
The conceptual impact matters most: the \emph{form} of the error penalty (linear and sample-weighted rather than squared) is what separates a tractable-but-slow method from an optimal one. Limitations: the assumptions (finite action sets, idealized actor updates, well-behaved weights) are far from deep-learning practice; the benchmarked algorithm departs from the analyzed one; and the elegant per-sample weighting collapses in practice to a single hand-searched constant per dataset.

\q{How were the results evaluated or validated?}
The theorems receive full appendix proofs built from standard statistical tools. The empirical side is thin: training curves only, a single baseline, no final-score table, no statistical tests, no ablation isolating the new penalty, and no numerical check of the theoretical rate improvement.

\q{What are the ethical implications? How are they addressed?}
No ethics discussion. Optimal-rate language could create false confidence for clinical or robotic use: the guarantees rest on assumptions that cannot be verified from logged data, concern average return rather than per-state safety, and do not cover the version actually benchmarked. Because the method competes only with behaviors the data cover, logging bias passes straight through. Public code is the one clear positive.

\q{What themes are suggested by this research? (minimum 3)}
\emph{Guarantees under partial data coverage}: optimal rates against covered behaviors only. \emph{Theory vs.\ practice}: the tested algorithm is not the analyzed one. \emph{Relative pessimism}: distrust improvement claims, not absolute values. \emph{Adaptive per-sample conservatism}: principled in theory, sacrificed for tractability. \emph{Sensitivity to settings the theory does not cover}.

\end{qna}
\tiein{A-Crab's per-sample reweighting of prediction errors is the theoretical cousin of BCA's calibrated widths: the paper proves such weighting attains the optimal statistical rate, and its practical collapse to one global constant marks exactly the gap BCA aims to fill.}

% ============================ PAPER 6 ============================
\paper{Paper 6: A Clean Slate for Offline RL (Unifloral)}
Rather than leading with yet another algorithm, this paper rebuilds twelve existing ones in one shared, fast codebase and grades them by a new yardstick: how well each scores for a given amount of live testing spent on tuning, a cost earlier papers rarely stated.
\begin{qna}

\q{What is the citation?}
Jackson, M. T., Berdica, U., Liesen, J., Whiteson, S., \& Foerster, J. N. (2025). A clean slate for offline reinforcement learning. \emph{Advances in Neural Information Processing Systems, 38}. Version reviewed: arXiv:2504.11453v1. \url{https://doi.org/10.48550/arXiv.2504.11453}

\q{What is the research area?}
Offline RL methodology, evaluation, and software rather than a new learning rule. Its problem family is hyperparameter and policy selection under limited online access, together with fair benchmarking. This places it beside earlier unification and reimplementation efforts (Rainbow, CORL) rather than beside algorithm papers.

\q{What is the research question?}
How should offline RL methods be defined, implemented, and compared so that reported scores honestly reflect the live interaction spent tuning them, and can one cleanly implemented design space, containing prior algorithms as settings, be used to discover better ones?

\q{What is the significance of the question?}
Offline RL's selling point, learning without touching the real world, is undermined if published scores quietly depend on hundreds of live test episodes used to pick settings. Numbers from different papers then mix algorithmic quality with tuning effort, and nobody can tell which is which; by 2025 the field had many competing methods, copied baselines, and incomparable codebases.

\q{What are the contributions reported in this paper?}
Four. A taxonomy of four offline RL settings, distinguished by when live testing is allowed, plus a definition of ``a method'' as an algorithm together with one fixed range of settings used on every dataset. An evaluation protocol in which a bandit (a strategy for choosing which trained policy to test next) simulates tuning, so performance is plotted against the live-testing budget. A family tree of prior methods with compact, fast reimplementations of twelve of them, and Unifloral, one configurable implementation covering most of the tree (it omits CQL). Finally, two genuinely new algorithms (TD3-AWR and MoBRAC) created purely by mixing settings, with no new code. All code is public.

\q{What is the research method?}
Conceptual analysis, protocol and framework construction, then benchmarking; no new theory. For each algorithm--dataset pair, 20 policies are trained with settings sampled from the original papers' ranges and their episode returns stored; 500 simulated tuning sessions then select among them under a stated budget, making the cost of selection visible. The study uses nine datasets, with timing runs on one GPU.

\q{What are the results?}
No prior method dominates: ReBRAC leads on five of nine datasets at some budget, IQL on four, and model-based methods never rise above sixth outside locomotion. The reimplementations are tiny (59 changed lines turn TD3+BC into ReBRAC) and dramatically faster (CQL trains in about 10 minutes vs.\ over three hours in a popular library). TD3-AWR beats each of its parent methods on most datasets (ReBRAC on six of nine, IQL on seven). And with only a few noisy test episodes, selection can be fooled by erratic ``distractor'' policies.

\q{What are the impacts and limitations of the contributions?}
The field gains a shared yardstick and toolkit: hidden tuning becomes an explicit, plotted cost, and component-level ideas can be tested by changing settings rather than writing new code. Limitations: only the pre-deployment setting is implemented; rankings depend on the original papers' tuning ranges; everything uses one benchmark suite; some fixed configuration choices degrade certain algorithms relative to their originals; and the two headline average speedups are not recoverable from the reported data.

\q{How were the results evaluated or validated?}
By stored-episode tuning simulations with confidence intervals, a reproduction table checking the reimplementations against published scores, timing runs, and a distractor-policy analysis over 100{,}000 orderings. Gaps: no formal statistical tests; the intervals reflect selection randomness over one pool of trained policies rather than retraining; some protocol details go unreported; and no check confirms the unified implementation matches each standalone file.

\q{What are the ethical implications? How are they addressed?}
There is no dedicated ethics section, but the paper is itself a research-integrity intervention: undisclosed live tuning inflates deployability claims, and pricing that budget makes the inflation visible. Unaddressed: selecting policies by average return ignores worst-case outcomes in the safety-critical domains the paper cites; distractor policies are exactly what a careless real selection process would deploy; and every benchmark is simulated. The efficiency gains widen access for smaller labs.

\q{What themes are suggested by this research? (minimum 3)}
\emph{Online tuning and offline model selection}: tuning is priced in episodes, and the selection rule itself matters. \emph{Evaluation rigor, reproducibility, fair comparison}: one codebase, one settings range per method, budget curves instead of headline numbers. \emph{Algorithms as configurations}: much of what separates published methods is settings, not principles. \emph{Benchmark dependence}. \emph{Behavior regularization}: the strongest methods still tie the policy to the data. \emph{Computational efficiency}: cheap training makes honest evaluation affordable.

\end{qna}
\tiein{Unifloral is the baseline reference bundled with BCA's code, and its budget-aware protocol is the right measuring stick: if calibrated per-sample caution reduces the search over global knobs, the payoff appears as better scores at small live-testing budgets.}

% ============================ PAPER 7 ============================
\paper{Paper 7: Quantile Q-Learning: Revisiting Offline Extreme Q-Learning with Quantile Regression (QQL)}
QQL replaces the most painful hand-tuned knob in a recent offline algorithm, a single ``optimism temperature'' per dataset, with a per-state scale learned from the data itself.
\begin{qna}

\q{What is the citation?}
Gao, X., Li, S., Cai, Y., \& Yu, W. (2026). Quantile Q-learning: Revisiting offline extreme Q-learning with quantile regression. \emph{Transactions on Machine Learning Research}. Version reviewed: arXiv:2511.11973v2. \url{https://doi.org/10.48550/arXiv.2511.11973}

\q{What is the research area?}
Offline RL, downstream of the in-sample family (IQL, XQL) in which values are learned only on recorded actions and the policy is extracted by weighted imitation. Its specific problem is temperature selection, or how optimistic the value estimates should be. It addresses this by borrowing quantile regression (learning specific percentiles rather than averages) from distributional RL.

\q{What is the research question?}
Can an XQL-style algorithm stay stable and competitive across many domains with one fixed configuration, if its temperature is estimated from the data as a per-state quantity instead of being tuned per dataset?

\q{What is the significance of the question?}
Per-dataset tuning normally means trying candidate policies live, requiring the interaction offline RL exists to avoid. XQL's temperature is the worst offender, with reproduced values spanning 0.5 to 10 across datasets. Because the temperature feeds both value learning and the policy's imitation weights, a bad choice destabilizes both at once. Making it a learnable per-state scale would remove the most sensitive knob.

\q{What are the contributions reported in this paper?}
A neat theoretical move: under a plausible noise model, the temperature equals a constant times the gap between two specific percentiles of the action-value distribution, so it can be learned by ordinary quantile regression instead of hand-tuning. Then QQL itself: two percentile value networks, a bias-corrected value target, and a policy update whose optimism automatically varies by state; a cautious generalization term for unseen actions; and a broad study (19 datasets from D4RL and NeoRL2, ablations, and checks of the noise model).

\q{What is the research method?}
Derivation under stated assumptions, then algorithm design and benchmarking. All settings are fixed everywhere (no per-dataset tuning) while baselines are allowed their tuned values; further studies sweep the two remaining knobs, test the noise model on synthetic data, and fit distributions to real prediction errors.

\q{What are the results?}
With one fixed configuration, QQL posts the best average on 10 of 17 standard datasets, and when XQL is forced to share one temperature too, QQL wins clearly (locomotion 85.6 vs.\ 73.6; mazes 84.9 vs.\ 67.3), with smoother training curves than its nearest relative. It is weak on both Hammer manipulation datasets and, notably, on the one safety-constrained task, where it scores below the dataset's own average. Compute is modest (about four hours per run, 1.6$\times$ its parent).

\q{What are the impacts and limitations of the contributions?}
QQL is evidence that XQL's most sensitive hyperparameter can be estimated from data, per state; the percentile-gap scale is a reusable idea. Limitations: the assumed noise model fits real errors only loosely (real errors are heavier-tailed); the two learned percentiles can cross, patched by hand; the learned scale carries no coverage guarantee; some ablation claims overstate the figures; one relevant baseline is missing; and no code link is given.

\q{How were the results evaluated or validated?}
By five-run averages (three-run on NeoRL2), ablations of both new components, knob sweeps, a synthetic test, and an error-distribution analysis. Weaknesses: no confidence intervals or statistical tests; baselines from mixed sources; an unjustified choice of shared temperature for the head-to-head comparison; a synthetic test that injects exactly the noise it then detects; and a few internally conflicting numbers.

\q{What are the ethical implications? How are they addressed?}
None discussed. The claim that fixed settings make the method practical for consequential use rests entirely on simulation, and the only safety-focused task is where it fails, with no reporting of constraint violations. A misestimated scale in sparsely covered states can mean \emph{less} caution exactly where mistakes cost most; logging bias is inherited. The fixed-configuration design does reduce risky live trial runs, which is the paper's implicit safety virtue.

\q{What themes are suggested by this research? (minimum 3)}
\emph{Hyperparameter sensitivity, relocated}: tuning moves from a global knob into a learned quantity with insensitive controls. \emph{Adaptive, state-dependent caution}: wider value spread across recorded actions triggers stronger caution. \emph{Uncertainty quantification without calibration}: a modeled scale with no coverage guarantee. \emph{Caution vs.\ generalization}; \emph{theory vs.\ practice}.

\end{qna}
\tiein{QQL is BCA's closest non-conformal precedent, using a per-state scale that acts on the same imitation weight that BCA modifies. However, its scale is learned under a parametric noise model with no guarantee; BCA's conformally calibrated width is the direct answer (though its coverage concerns held-out residuals, not the policy), and QQL's own heavy-tailed errors argue for it.}

% ============================ PAPER 8 ============================
\paper{Paper 8: UNIQ: Conformal Calibration for Adaptive Conservatism in Offline Reinforcement Learning (UNIQ)}
UNIQ gives a standard offline learner (IQL) a per-state caution dial driven by calibrated ensemble disagreement, with more caution where data are sparse and less where they are dense. It does so at a fraction of the memory cost of large-ensemble methods.
\begin{qna}

\q{What is the citation?}
Upadhyay, A. (2026, July 11). \emph{UNIQ: Conformal calibration for adaptive conservatism in offline reinforcement learning} [Poster presentation]. ICML 2026 Workshop on Decision-Making from Offline Datasets to Online Adaptation: Black-Box Optimization to Reinforcement Learning, Seoul, South Korea. Version reviewed: arXiv:2606.07592v1. \url{https://doi.org/10.48550/arXiv.2606.07592}

\q{What is the research area?}
Offline RL, built on IQL. The problem family is pessimism under distribution shift: deciding how much to distrust value estimates in poorly covered situations. Unlike global-knob methods (CQL, IQL, TD3+BC), huge-ensemble methods (SAC-N, EDAC), or learned-model penalties (MOPO), UNIQ pairs a small ensemble of value heads with conformal calibration, a statistics recipe that turns held-out prediction errors into a unit-free uncertainty scale.

\q{What is the research question?}
Can IQL's single global optimism setting be replaced by a per-state setting driven by conformally normalized ensemble disagreement, with tighter caution where data are sparse and looser caution where data are dense, without the memory cost of large ensembles?

\q{What is the significance of the question?}
The best setting of every mainstream caution knob changes per dataset and is normally found by trying policies live, an option unavailable offline. Raw ensemble disagreement, the usual uncertainty signal, has no common scale across tasks, so fixed penalties transfer poorly. A calibrated, lightweight, per-state rule would cut both live tuning and hardware demands at once.

\q{What are the contributions reported in this paper?}
A three-part mechanism: a small ensemble of value heads trained at three optimism levels, whose spread estimates uncertainty; a conformal step that measures prediction errors on held-out data and converts raw spread into a unit-free ratio; and a controller mapping that ratio to a state-specific optimism setting, with a matching penalty in the value targets. Supporting contributions: a formal appendix (the coverage property, controller behavior, IQL as a special case), a candid benchmark comparison that keeps the losing rows, a memory analysis, and ablations.

\q{What is the research method?}
Algorithm design with light formal analysis and benchmarking. UNIQ uses three heads per optimism level; every 5{,}000 training steps, held-out transitions recalibrate the error scale; the normalized uncertainty drives the per-state setting between 0.5 and 0.90/0.95. Two hand-set configurations cover nine simulated locomotion datasets over three seeds, with baselines copied from a 2023 benchmark report rather than rerun.

\q{What are the results?}
UNIQ averages 85.2 across the nine tasks versus IQL's 81.6, beating IQL on every task (by up to 8.1 points) and ranking third overall behind two much heavier or more heavily tuned methods. Its measured memory use is 250 MB against an estimated 2.5 GB for the ensemble leader. The ablations temper the story: a hand-set penalty coefficient matters more than either calibration or adaptivity, and a single head per level (so no uncertainty signal) nearly matches the full method (59.1 vs.\ 59.4).

\q{What are the impacts and limitations of the contributions?}
UNIQ is a proof of concept that conformal calibration can be added cheaply to a standard learner, with unusually honest reporting of losses. Limits: the conformal guarantee covers a task-wide error scale, not individual states; it holds only at each recalibration of a moving ensemble; the gains are partly confounded with per-task penalty settings; the evidence is locomotion-only; and several numbers conflict across tables, text, and appendix.

\q{How were the results evaluated or validated?}
By benchmark scores over three seeds against eight baselines copied (not rerun) from a prior report, single-seed ablations with inconsistent labeling, and analytic rather than measured speed and baseline memory estimates. Realized coverage and the distribution of the uncertainty signal are never reported, so the headline 3.6-point gain over IQL is plausible but weakly established, and the IQL comparison is not matched.

\q{What are the ethical implications? How are they addressed?}
No dedicated treatment; a frank scope section substitutes. The language of distribution-free finite-sample guarantees could imply a safety assurance the method does not provide: the coverage concerns prediction errors on logged data, not the states the learned policy will visit, and the controller grows \emph{more} optimistic wherever ensemble heads agree, which is dangerous if they share the same error. Data bias is inherited; validation is simulated; no code release is mentioned; the small memory footprint widens access.

\q{What themes are suggested by this research? (minimum 3)}
\emph{Adaptive, state-dependent caution}: one global setting is too careful in dense regions, too lax in sparse ones. \emph{Uncertainty quantification and calibration}: raw disagreement needs a common scale. \emph{Conformal guarantees and their limits}: coverage on residuals does not transfer to the derived signal or to visited states. \emph{Sensitivity migrates}: a new knob quietly inherits the old knob's tuning burden. \emph{Computational efficiency}; \emph{evaluation rigor}.

\end{qna}
\tiein{UNIQ is BCA's closest prior work: both bring conformal ideas to prediction-error-driven caution. UNIQ calibrates the errors but then normalizes ensemble \emph{spread} with that quantile, so the guarantee never reaches its uncertainty signal; BCA calibrates the errors against a fitted scale model with a WBCP threshold, so width and signal describe the same quantity, with the calibration frozen between periodic refreshes and portable across four host algorithms.}

\section*{Part B: Themes and Observations}

Two meta-themes dominate all eight papers. The technical one is distribution shift. Value estimates for actions missing from the data are the shared enemy, and every algorithm offers a different way of dosing pessimism: penalizing critics (CQL), holding policies near the data (TD3+BC, ReBRAC), learning values almost entirely from recorded actions (IQL, QQL, UNIQ), or using provable pessimism (A-Crab). The methodological meta-theme is nearly as common and encompasses hidden engineering, hidden live tuning, and evaluation rigor. TD3+BC, ReBRAC, and Unifloral collectively show that reported progress often reflects engineering choices and tuning budgets rather than stated ideas. Less prevalent but most significant is adaptive, state-dependent caution with calibrated uncertainty (A-Crab's weights, QQL's learned scale, UNIQ's conformal ratio). It attacks the shared root weakness of one global caution knob tuned per dataset and directly motivates BCA's conformally calibrated, per-sample caution.

\end{document}
